\begin{abstract}
Financial institutions need intrusion detectors that are accurate and
auditable and that do not themselves become a channel for leaking customer
data. We present HALX, a framework for proactive threat detection in
financial networks with three parts: a feature-token Attention-LSTM
encoder, an XGBoost scoring head, and a privacy-preserving inference layer.
The privacy layer enforces Mondrian $k$-anonymity on quasi-identifiers and
releases the preprocessing statistics under $(\varepsilon,\delta)$-differential
privacy. Every alert carries a threat score, a severity tier and the
most-attended features as an explanation. We compare HALX with seven
baselines over three seeds on the official UNSW-NB15 split and on a
CICIDS2017 subset, giving every method the same privatized inputs. On
CICIDS2017, HALX reaches an F1-score of \cicHALXfOne\% at a
\cicHALXfpr\% false-positive rate. This is on par with the strongest tree
ensembles and gives 2--10$\times$ fewer false alarms than CNN-LSTM, BiLSTM
and Transformer detectors. On UNSW-NB15, a train--test distribution shift
holds every model near 89\% F1 on average. There HALX matches XGBoost
(\unswHALXfOne\%), has the most stable false-positive rate across seeds,
and detects the most attacks at a 1\% false-alarm budget
(\unswHALXtprOne\%, seed 0). Under this shift,
\unswsevhigh\% of its high-severity alerts are real attacks. The privacy
layer ($k{=}10$, $\varepsilon{=}1$) costs at most \PRIVCOSTMAX{} F1 points,
and HALX scores thousands of flows per second on a CPU. Code and results
are publicly available.
\end{abstract}
